\section{人身犯罪 - 侵犯生命、身体的犯罪}
\begin{enumerate}
  \item \zd{保护法益 \yjt 构成要件（主体、行为、对象）\yjt 法定刑升格条件 \yjt 犯罪形态 \yjt 共同犯罪 \yjt 罪数}。
  \item 罪名大多动宾结构。动词 + 名词。
  \item 故意杀人。他人的生命权。
        \begin{enumerate}
          \item 安乐死。承诺放弃生命是无效的，可以从宽。
          \item (欺骗 \hz 强迫)自杀、自伤、自残 \yjt 间接正犯。\al{相约自杀，自己逃走；欺骗精神病患者自杀、自残}。
          \item 教唆或帮助他人。
                \begin{enumerate}
                  \item 自杀。观点展示。
                        \begin{enumerate}
                          \item \zd{有罪}。生命要绝对保护（身体健康不用），共犯从属性可以有例外。
                          \item 无罪。共犯从属性。
                        \end{enumerate}
                  \item 自伤、自残。无罪。\zd{实行者是本人}。共犯从属性；没必要绝对保护。
                  \item 被害人承诺。承诺放弃生命、重伤无效。实行者不是本人。
                \end{enumerate}
        \end{enumerate}
  \item 故意伤害。他人的身体健康。
        \begin{enumerate}
          \item 致人\zd{轻伤} \yjt\, 16 岁；致人重伤 \hz 死亡 \yjt\, 12 岁。
          \item 伤害行为有程度。
                \begin{enumerate}
                  \item \zd{轻微（刑法不管）}、轻伤、重伤、伤害致死。
                  \item 一般是身体伤害，精神伤害具体分析。\al{假牙、假肢、隐形眼镜具体分析}。
                \end{enumerate}
          \item 伤害行为有要求。\al{推一把}。
                \begin{enumerate}
                  \item 主观上要有造成\zd{轻伤的意图}。
                  \item 客观上行为要有\zd{造成轻伤的一般可能性}。
                \end{enumerate}
        \end{enumerate}
  \item 对死亡态度不一样。
        \begin{enumerate}
          \item 故意杀人 \yjt 希望 \hz 放任。
          \item 故意伤害致人死亡 \yjt 过失心理。
          \item 过失致人死亡 \yjt 缺少伤害故意。
        \end{enumerate}
  \item 故意杀人、故意伤害。性质相同、程度不同。包容评价。\al{伤害 2 刀，杀害 2 刀，只有 1 刀致命}。
  \item 组织出卖人体器官罪。打击器官买卖。
        \begin{enumerate}
          \item 活体器官，不包括尸体器官、血液（非法组织卖血罪），\zd{看能不能移植}。
          \item 出卖者不构成此罪。
          \item \zd{出卖行为必须有效的承诺}。\yi 年满 18 \er 是否受到欺骗、胁迫 \san 非法买卖还是合法捐赠。
        \end{enumerate}
  \item 遗弃罪。不作为犯。
        \begin{enumerate}
          \item 和不作为的故意杀人罪区别：\zd{生命危险是否紧迫}。\al{扔到马路边/深山老林}。\zd{性质相同，程度之分}。
          \item 不限于家庭成员，抚养义务要扩大解释。
        \end{enumerate}
\end{enumerate}
